% Cortes mecánicos íntegros: parte_iv.tex:55--106,337--408.
% Residencia material del ensamblaje sucesor de 102 capítulos.
% Residencia de realización: capítulo 76.
\section{Spatial suspension}

Let
\[
 \Gamma(\theta)=
 \bigl(r(\theta)\cos\theta,r(\theta)\sin\theta,z(\theta)\bigr).
\]
Its curvature and torsion are calculated by means of
\[
 \kappa_{\mathrm{F}}=\frac{\|\Gamma'\times\Gamma''\|}{\|\Gamma'\|^3},
 \qquad
 \mathcal T_{\mathrm{F}}=
 \frac{\det(\Gamma',\Gamma'',\Gamma''')}
 {\|\Gamma'\times\Gamma''\|^2}.
\]
A general helix in the sense of Lancret satisfies
\[
 \frac{\mathcal T_{\mathrm{F}}}{\kappa_{\mathrm{F}}}=\text{constant}.
\]
An arbitrary radial suspension need not satisfy that condition. The
name \emph{radial helical suspension} preserves the winding structure
without asserting an unproved differential property.

\section{Geometric trichromy}

The generators
\[
 J_{+1}^2=-I,\qquad J_0^2=0,\qquad J_{-1}^2=I
\]
produce elliptic, parabolic, and hyperbolic evolutions. A history may
alternate them and publish a curve whose Frenet curvature varies
continuously. The \emph{trichromy} describes the sequence of internal regimes;
Frenet geometry describes the output. The radial reader preserves both
without confusing them.

\begin{figure}[htbp]
\centering
\begin{tikzpicture}[>=Latex,node distance=10mm and 16mm,
 n/.style={draw=HMTNavy,fill=white,text width=34mm,minimum height=10mm,align=center,font=\small\sffamily}]
\node[n] (int) {historia TRIT\\\(+1,0,-1,\ldots\)};
\node[n,right=of int] (rad) {reader radial\\\(p,\Lambda,C,\mu\)};
\node[n,right=of rad] (curve) {published curve\\\(r(\theta),z(\theta)\)};
\node[n,below=of curve] (frenet) {Frenet geometry\\\(\kappa_{\mathrm{F}},\mathcal T_{\mathrm{F}}\)};
\draw[->,HMTBlue,thick] (int)--(rad);
\draw[->,HMTBlue,thick] (rad)--(curve);
\draw[->,HMTBlue,thick] (curve)--(frenet);
\draw[->,HMTRed,dashed,bend left=28] (frenet.west) to node[below,font=\scriptsize\sffamily]{does not select retrospectively} (int.south);
\end{tikzpicture}
\caption{Causal separation between internal regime and visible curvature.}
\label{espiral:fig:trit-frenet}
\end{figure}

\chapter{Dynamic structure of space and consequences}

\section{Helicoidal leaf lamination}

The geometric synthesis is a lamination whose elements are lifted
histories. Each leaf carries:
\begin{enumerate}
\item an angular phase;
\item a coordinate radial internal;
\item a memory depth;
\item a sequence of sheets APP;
\item a local TRIT regime;
\item a TPK holonomy and its boundary.
\end{enumerate}
Transverse sections are circles or level curves; families of
sections form cylinders; radial variation produces suspensions; the
involution produces conjugate sheets. Real geometry is the image of this
lamination under a publication that may contract several sheets.

\section{Kinematic classification by angular and radial increments}

A radial line appears when \(\omega=0\) and \(\beta\neq0\). A
circle appears when \(\omega\neq0\) and \(\beta=0\). The spiral
appears when both are nonzero. These three outputs are configurations of
the same transport mechanism, not unrelated objects. The APP support
is two-dimensional from the outset; the one-dimensional traversal is a
particular section of that support.

\section{Mobius Quotient via Return with Orientation Inversion}

When the route involution is identified at the seam after one
turn, the quotient may acquire a Möbius band. The assertion requires
a free action and the intertwiner that transports orientation. The band does not
follow from the mere presence of a sign \(-1\); it follows from the seam
\[
 (x,0)\sim(C_{\mathrm{rev}}x,1)
\]
in the declared domain. In the material atlas, the dodecaphasic involution and
the cellular inversion remain distinct operations.

\Needspace{34\baselineskip}
\section{Classification of 3+2 equilibrium points via the Hessian}

The intuition of Lagrange's three collinear and two triangular points may be
formulated as a subsequent proof. Given an effective potential
\(V_{\mathrm{eff}}\) produced by the HMT--MD realization, the critical points
satisfy
\[
 \nabla V_{\mathrm{eff}}=0.
\]
Its classification requires the Hessian and the Morse index. Only if the
dimensional map produces exactly three collinear and two triangular critical points
can it be related to the partition \(3+2\). The correspondence is not used
as an input to the spiral taxonomy.

\section{Causal hierarchy between curves, fields, and material sections}

The spiral correspondence opens three coordinated outputs:
\[
 \text{APP--TRIT--TPK ledger}
 \longrightarrow
 \begin{cases}
 \text{Archimedean curve},\\
 \text{electromagnetic mode},\\
 \text{material route signature}.
 \end{cases}
\]
The first is algebraically closed in this work. The second has
an exact realization for circular modes and a spectral
compatibility condition. The third contains the central electron and the
atlas but requires verification of the constancy square on each
multisection before assigning flavor taxonomies.
